\section{Conclusion}
\label{sec:fazit}
The investigation into prompt injection defenses has transitioned from basic heuristic filters to sophisticated system-level frameworks and preference-optimized models.
This final chapter summarizes the findings to answer the core research question regarding the effectiveness of instruction-data separation.
\subsection{Summary of Structural Vulnerability}
Prompt injection is a structural vulnerability inherent to any architecture that treats control instructions and user input as a single, concatenated token sequence.
Because modern LLMs process natural language as a medium that simultaneously serves as both data and executable instruction, they lack the native proficiency required to recognize malicious intent.\\
Furthermore, an inverse scaling effect often persists: as models become more intelligent and better at following complex instructions, they simultaneously become more effective at executing malicious payloads if the initial security boundary is bypassed.
\subsection{Semantic vs. Architectural Separation}
The research highlights a fundamental divide between semantic and architectural separation.
Semantic separation, achieved via fine-tuning techniques like StruQ and SecAlign, significantly raises the cost of attacks and effectively blocks most heuristic manual injections.
However, these defenses remain probabilistic and brittle against adaptive attacks, because they rely on learned behavior that can be subverted by architecture-aware optimization.\\
In contrast, architectural separation through systems like f-secure and CaMeL offers formal security guarantees based on Information Flow Control.
These systems treat the LLM as a black box and physically restrict the flow of untrusted data to the "planner," providing robustness that persists even as underlying models are upgraded.
\subsection{The Adaptive and the Security "Arms Race"}
The evolution of prompt injection constitutes an ongoing "arms race".\\
Heuristic defenses are bypassed by optimization-based attacks like GCG.
As fine-tuning defenses improved, new architecture-aware attacks like ASTRA emerged to target the internal attention matrices of transformers, breaking even the strongest preference-optimized models.
These findings suggest that robustness is often a sign of a lack of strong optimization signals in tests, rather than a full remediation of the underlying vulnerability.
\subsection{Future Outlook: Defense in Depth}
Ultimately, there is no "silver bullet" for prompt injection.
The path forward lies in defense in depth, combining model-level robustness with architectural isolation, least-privilege tool access, and human-in-the-loop approvals for sensitive operations.\\
Future research must address the emerging challenges of multi-modal injections and increasingly autonomous agents, which expand the attack surface beyond traditional text-based interfaces.\\\\
As Large Language Models become the backbone of autonomous digital systems, moving beyond a "unsafe-by-design" architecture is the ultimate challenge for the next generation of AI security. Establishing verifiable security guarantees while concurrently maximizing utility will be the prerequisite for the safe deployment of LLMs in load-bearing software infrastructure.